\begin{center}
\small
\begin{tabular}{@{}lrrrr@{}}
\toprule
Unrestricted integer presentation & Variables & Equations & Operations & Degree\\
\midrule
Root-eliminated product & 32 & 1 & 138 & 584\\
Normalized compact product & 33 & 1 & 130 & 304\\
Ordinary double-gap product & 33 & 1 & 134 & 159\\
Ordinary quadratic lift with positivity relations & 60 & 33 & --- & 2\\
Its sum of squares & 60 & 1 & --- & 4\\
\bottomrule
\end{tabular}
\end{center}
The ternary positivity prefix uses seven operations per positive port, hence adds 63 operations to each compact circuit. The 60-variable system instead retains the 33 original coordinates and adds 27 variables and nine quadratic relations $v=\PP(u_v,v_v,r_v)$.

The degree counts use weights two for the nine positive ports and one for all remaining coordinates, before substitution. The factor weighted degrees of the 14-variable product are $44,35,13,4,488$, summing to 584. For the ordinary double-gap product they are $24,35,39,13,4,44$, summing to 159. The latter is not twice degree 83: its gap roots remain signed and have weight one. The weighted leading forms are nonzero at the actual $\beta>0$. Disjoint positivity blocks preserve this fact, since setting each block's last two coordinates to zero leaves the injective substitution $v\mapsto u_v^2$ on leading forms. Thus these are exact degrees, not only substitution upper bounds.

Further shipped tradeoffs include 47 unrestricted variables at degree eight, 43 at degree ten, and 38 at degree sixteen. Their literal residuals and independent receipts are supplied, but the table highlights the clearest choices rather than treating every elimination as a separate main result.
